\documentclass[a4paper,11pt]{article}
\pagestyle{empty} % No page numbers
\usepackage{geometry}
\geometry{margin=.7in}
\usepackage{enumitem}
\usepackage{hyperref}
\usepackage{xcolor}
\usepackage{tabularx}
\usepackage[normalem]{ulem}

\def\jobh{\hspace{-1ex}}
\def\locationh{\hspace{7cm}}
\def\dateh{\hspace{3cm}}
\def\titlev{\vspace{.5cm}}
\def\sectionv{\vspace{.2cm}}
\def\subsectionv{\vspace{-5.3ex}}
\def\betweenpoints{.1cm}
\def\postpointsv{\vspace{0cm}}
\def\betweenjobsv{\vspace{0.4cm}}
\def\betweenskillsh{\hspace{1.4cm}}


\hypersetup{
    colorlinks=true,
    urlcolor=blue,
}

\begin{document}

\begin{center}
    {\LARGE \textbf{Carlos Serrouya}}\\[0.3cm]
    Calgary \hspace{0.5em} 
    \vrule width 0.4pt height 1.2em depth 0.2em \hspace{0.5em}
    \href{https://www.linkedin.com/in/carlos-serrouya-27053b276}{LinkedIn} \hspace{0.5em} 
    \vrule width 0.4pt height 1.2em depth 0.2em \hspace{0.5em}
    \href{mailto:carlos.serrouya@ucalgary.ca}{carlos.serrouya@ucalgary.ca} \hspace{0.5em} 
    \vrule width 0.4pt height 1.2em depth 0.2em \hspace{0.5em}
    250-837-1041
\end{center}
\titlev

\noindent \makebox[\textwidth][c]{\textsc{Work Experience and Personal Projects}}
\vspace{-2.35em} % Adjust this value as needed
\par\noindent
\rule{\textwidth}{0.4pt}

\sectionv

\noindent \begin{tabbing}
    \locationh \= \kill
    \textbf{Data Science Intern} \textit{Open Ocean Robotics} \` 2025 – Present (1 Year) \\
\end{tabbing}

\begin{itemize}[leftmargin=.5cm, itemsep=.1cm, before=\subsectionv, after=\postpointsv]
    \item Key contributor in taking \textbf{Enhanced Horizon\texttrademark} from development to a robust object tracking and localization system used on autonomous surface vessels.
    \item Deployed ML applications on edge devices and cloud using Docker, PyInstaller, and Git CI pipelines; optimized latency and throughput with TensorRT, NumPy, and multithreading
    \item Implemented an end-to-end model benchmarking pipeline that auto-generates reports and visualizations for each new model release, enhancing reproducibility and performance tracking
    \item Improved object detection models developed with PyTorch by +20\% inference speed; mAP50: 28→45, mAP50-95: 7→20, precision: 40→80\%, recall: 30→40\%, class confidence: 80→90\%
    \item Developed distance estimation object tracking algorithms using camera ray tracing, Kalman filter, statistical propegation of errors; deployed from simulation to edge devices on USVs with highly positive field feedback from USV operators, earning a company wide innovation award
    \item Automated large-scale data uploading and labeling pipelines, processing terabytes every two weeks, reducing manual upload time by ~3 hours per cycle
    \item Worked closely with a fellow intern during Python-to-Rust migration of key DS applications, explaining statistical and mathematical concepts and assisting with debugging
    \item Conducted advanced and extensive statistical analysis of data science applications performance; visualized results using matplotlib, seaborn, grafana, and \textbf{Xploreview\texttrademark} to guide engineering improvements and inform stakeholders. Presented as Visual lead to multiple clients totalling in hundreds of thousands in contracts.
\end{itemize}


\noindent \begin{tabbing}
    \locationh \= \kill
    \textbf{Data Science Intern} \textit{Alberta Biodiversity Monitoring Institute} \` May-September 2023, 2024\\
\end{tabbing}


\begin{itemize}[leftmargin=.5cm, itemsep=.1cm, before=\subsectionv, after=\postpointsv]
    \item Developed a proof-of-concept application to estimate snow coverage and depth from camera trap images using a fixed-reference-pole semantic segmentation and classification approach; done by using Tensor flow for U-net model and transfer learning to random forest model for classification model. Was able to reliably estimate snow depth in a given image within 7cm
    \item Created object contouring, detection, and outlier detection algorithms, to indetify inconsisensies in application results
    \item Managed large datasets with Pandas and SQL, integrated basic statistical models as SQL triggers.
    \item Presented machine learning results to a non-technical panel, achieving 84\% snow depth accuracy and 96\% coverage classification.
    \item Reduced false positives in Microsoft's Megadetector by 72\%, false negatives by 9\%, by implementing multiple regression post process model
    \item Visualized model performance data with Matplotlib and Seaborn.
\end{itemize}

\betweenjobsv

\noindent \makebox[\textwidth][c]{\textsc{Education and Certifications}}
\vspace{-2.35em} % Adjust this value as needed
\par\noindent
\rule{\textwidth}{0.4pt}

\sectionv

\noindent
\begin{minipage}[t]{0.45\textwidth} % Left column (45% of text width)
    \textbf{University of Calgary} \\
    \textit{BSc Computer Science} \hfill 2022-2027 \\
    \textit{BSc Mathematics} \hfill 2021-2026
\end{minipage}
\hfill % This creates space between the columns
\begin{minipage}[t]{0.45\textwidth} % Right column (45% of text width)
    \textbf{Coursera} \\
    \textit{Advanced Learning Algorithms} \hfill 2023 \\
    \textit{Regression \& Classification} \hfill 2023
\end{minipage}

\betweenjobsv
\sectionv

\noindent \makebox[\textwidth][c]{\textsc{Advanced Skills}}
\vspace{-2.35em} % Adjust this value as needed
\par\noindent
\rule{\textwidth}{0.4pt}

\sectionv

\noindent
\begin{minipage}[t]{0.32\textwidth} % Left column (32% of text width)
    \textbf{Microsoft Office} \\
    \textbf{Testing Automation} \\
\end{minipage}
\betweenskillsh
\begin{minipage}[t]{0.32\textwidth} % Middle column (32% of text width)
    \textbf{Data Validation} \\
    \textbf{Data Visualization} \\
\end{minipage}
\betweenskillsh
\begin{minipage}[t]{0.32\textwidth} % Right column (32% of text width)
    \textbf{SQL} \\
    \textbf{Power BI} \\
\end{minipage}

\end{document}



